\chapter{Scaling and resource limits}
\label{dnad24:chapter}
\ChapterGoals{%
\item freeze a resource vector before comparing designs, and refuse single-number comparisons;
\item extend the sampling bound with Poisson coverage and convert probability into matter and volume;
\item explain how diversity dilutes encounters at fixed total concentration, and why depth sets latency;
\item estimate diffusion times and the probability that a correct molecule reaches the observed sample;
\item state what an energy claim must include, from Landauer's bound to laboratory overheads.}
{Chapters~4, 7 and~23.}
A molecular program can have a short description and require a large physical
population. A reaction can occur in parallel and still take a long time. A tube
can contain the correct answer and produce an unreadable observation. These are
three different failures of scale.

Chapter 4 introduced multiplicity, diversity and witness recovery. Chapter 23
showed that local reliability does not determine whole-object correctness.
We now combine those ideas into a resource analysis that exposes where a design
moves its cost. The new tool is not another claim that DNA is massively parallel.
It is an executable account of what must grow when the input grows.

\section{Begin with a deliberately naive search}
Suppose an $n$-bit search represents each of its $M=2^n$ candidate strings by
a distinct molecular identity. A filter removes invalid candidates; a reader
returns a surviving witness. Assume for this example that all candidates are
generated uniformly and that the representation itself is correct.
What must exist before the filter can exploit its parallelism?

At least one instance of each required identity must be present if the claim is
complete enumeration. This requirement is stronger than having $M$ molecules:
some identities can occur repeatedly while others are absent. Conversely,
finding one witness need not require enumerating every candidate. A tailored
algorithm can change the representation, generation distribution and required
coverage. Our calculation is a cost of this architecture, not a lower bound on
all molecular algorithms for the underlying search problem.

Early DNA-computing research supplied important algorithms for molecular
satisfiability and related problems \cite{Boneh1996}. Here we do not equate
a small number of high-level operations with a small physical resource budget.
A separation operation acts on a population whose preparation and observation
must also be accounted for.

\Plate{01-resource}{Parallel work is only one component of the resource path.
Candidate inventory, serial reaction dependencies and finite readout capacity
remain separate costs. The graph is an accounting model, not a laboratory protocol.}
{dnad24:resource}

\section{Freeze the resource vector before comparing designs}
Write a design's resource account as
\begin{equation}
\mathcal R(n)=
\bigl(S,N_{\mathrm{mol}},V,c_{\mathrm{tot}},L_{\mathrm{nt}},
d_{\mathrm{dep}},t_{\mathrm{run}},Q_{\mathrm{read}},E_{\mathrm{used}}\bigr).
\label{dnad24:vector}
\end{equation}
Here $S$ is the number of chemically distinct species; $N_{\mathrm{mol}}$ counts
physical molecules; $V$ is volume in litres; $c_{\mathrm{tot}}$ is total molar
concentration; $L_{\mathrm{nt}}$ is total synthesized nucleotide inventory;
$d_{\mathrm{dep}}$ is serial dependency depth; $t_{\mathrm{run}}$ is a declared
latency; $Q_{\mathrm{read}}$ is the amount of readout demanded; and
$E_{\mathrm{used}}$ is actual energy expenditure when measured.

Specify when each quantity is counted. Peak molecules during amplification differ
from input molecules. Number of reusable tile types differs from number of copies
incorporated in an assembly. A seed prepared separately is not a free initial
condition. Chapter 22's four computational XOR tiles did not describe the full
boundary-seed inventory.

A resource vector also requires a success contract: which event, with what
probability, by what deadline? An answer to a decision problem can need only one
distinguishable bit; reconstructing every candidate string demands much more.
No scalar operations-per-second number summarizes all of these designs.

\section{Poisson coverage: an exact extension of the sampling model}
Chapter 4 counted fixed draws and bounded missing identities. Here we use a
different, explicitly \Term{Poissonized} preparation model. For each identity
$j=1,\ldots,M$, let the copy count $C_j$ be independent Poisson with mean $\mu$.
This describes random total count, not a tube with exactly $M\mu$ molecules.
Independence and equal abundance are assumptions, not consequences of using DNA.

Since $P(C_j=0)=e^{-\mu}$, complete representation has probability
\begin{equation}
P_{\mathrm{cover}}=(1-e^{-\mu})^M.
\label{dnad24:coverage}
\end{equation}
To make the probability of missing any required identity at most
$\delta\in(0,1)$, solve rather than guess:
\begin{equation}
\mu_{\mathrm{exact}}
=-\log\!\left[1-(1-\delta)^{1/M}\right].
\label{dnad24:mu}
\end{equation}
The expression inside brackets can be tiny. Evaluating two nearly equal numbers
and subtracting them loses significant digits. The implementation instead uses
\[
\mu_{\mathrm{exact}}
=-\log\!\left[-\operatorname{expm1}
  \left(\frac{\operatorname{log1p}(-\delta)}{M}\right)\right].
\]
\Listing{coverage_copies}

The union bound gives a simpler sufficient condition:
\[
P(\exists j:C_j=0)\le M e^{-\mu},
\qquad \mu\ge\log M+\log(1/\delta).
\]
It does not require missing events to be independent once their marginal
probabilities are known; the exact product does. At fixed $\delta$ and large $M$,
Equation~\ref{dnad24:mu} gives
$\mu=\log M-\log[-\log(1-\delta)]+o(1)$.
Thus full representation requires expected total inventory of order
$M\log M$ in this particular random-preparation model, not just $M$.

\Plate{02-coverage}{Copy count and species coverage answer different questions.
The plate shows eight identity classes and an independent-Poisson occupancy model,
not eight physically resolved molecular species in an actual solution.}
{dnad24:occupancy}

\subsection{A falsifier for the coverage assumption}
If one required identity has generation probability zero, no increase in the
copies of other identities can restore complete coverage. If abundances vary,
replace the equal-mean expression by
$\prod_j(1-e^{-\mu_j})$ under independent Poisson counts. A rare species can then
dominate failure even when total DNA is abundant.

With fixed total draw count, occupancy events are dependent and
Equation~\ref{dnad24:coverage} is no longer exact. The reference includes a small
dynamic-programming oracle: after $k$ uniform draws, let $p_k(s)$ be the probability
that $s$ distinct identities have been seen. A next draw preserves $s$ with
probability $s/M$ and increases it with probability $(M-s)/M$.
The routine updates those two transitions for each state.
\Listing{fixed_draw_complete}
Exhaustive enumeration for two and three species independently verifies the
recurrence. A separate exponential-arrival simulation samples a Poisson total and
checks the product model. Agreement between these oracles is meaningful only when
their preparation contracts are the same.

\section{Translate the probability into matter}
The SI mole fixes $N_A=6.02214076\times10^{23}\ \mathrm{mol}^{-1}$ exactly
\cite{BIPMmole}. For molar concentration $c_{\mathrm{tot}}$ and volume $V$ litres,
\[
\mathbb E[N_{\mathrm{mol}}]=N_Ac_{\mathrm{tot}}V.
\]
A uniform pool with mean $\mu$ copies of each of $M$ identities consequently needs
\begin{equation}
V_{\mathrm{model}}=\frac{M\mu}{N_Ac_{\mathrm{tot}}}.
\label{dnad24:volume}
\end{equation}
\Listing{coverage_volume}
\begin{center}\small\input{\Generated/scaling.tex}\end{center}
The generated table assumes complete coverage with probability at least .95,
a uniform independent-Poisson pool, and total concentration $1\ \mu\mathrm M$.
The values are computed requirements under those assumptions, not feasible
synthesis orders, experimental volumes or universal complexity lower bounds.

All candidates are charged to the same total concentration. Assigning
$1\ \mu\mathrm M$ to each of $M$ species instead would raise total concentration
to $M\ \mu\mathrm M$ and silently change the physical problem. Increasing $n$ by
one doubles $M$ and also increases the required copies per identity. Under the
model, fixed-concentration volume grows slightly faster than pure doubling.

If each candidate has $\ell$ nucleotides, expected candidate nucleotide inventory
is $\ell M\mu$. Primers, gates, fuels, reporter strands and control species add
their own inventories. Turning this into mass requires the actual molecular
composition and modification state, not an unexplained universal molecular weight.
A computational representation that shares prefixes or uses assembly rather than
whole-string enumeration can change the inventory law; that is an algorithmic
change to be analyzed, not a reason to omit material accounting.

\section{Diversity dilutes encounters at fixed total concentration}
Consider one target associating with one required species in a well-mixed pool.
With association constant $k_{\mathrm{on}}$ in
$\mathrm{M}^{-1}\mathrm{s}^{-1}$ and species concentration $c_j$ in M,
the pseudo-first-order encounter rate is $r_j=k_{\mathrm{on}}c_j$.
For an exponential waiting-time model,
\[
\mathbb E[T_j]=1/r_j.
\]
Under equal concentrations $c_j=c_{\mathrm{tot}}/S$,
\begin{equation}
\mathbb E[T_j]=\frac{S}{k_{\mathrm{on}}c_{\mathrm{tot}}}.
\label{dnad24:wait}
\end{equation}
Increasing diversity at fixed total concentration can therefore slow a particular
required encounter. It does not necessarily slow every reaction by the same factor:
different pathways have different stoichiometry, concentrations and competition.

The generated results include an internal field evaluating this formula at assigned
$k_{\mathrm{on}}=10^6\ \mathrm{M}^{-1}\mathrm{s}^{-1}$; that rate is a teaching
parameter, not a measured value for the encoded species. Irreversibility,
constant partner concentration and well-mixed transport are assumptions. If a
partner is depleted or a contact dissociates, the one-clock model is insufficient.

\Plate{03-concentration}{At fixed total inventory and volume, more distinct
species means fewer copies and lower concentration per species. Encounter arrows
refer to an assigned required partner, not every molecular collision.}
{dnad24:concentration}

\section{Parallelism does not erase dependency depth}
Let a successful computation require $d$ serial stages. In an ideal irreversible
chain, assume each stage has an independent exponential waiting time with common
rate $r>0$ and that the next stage cannot start before its predecessor completes.
Then
\[
T_{\mathrm{chain}}=\sum_{\ell=1}^{d}T_\ell,\qquad
\mathbb E[T_{\mathrm{chain}}]=d/r.
\]
For integer $d\ge1$, the completion probability is the Erlang CDF
\begin{equation}
F_d(t)=1-e^{-rt}\sum_{k=0}^{d-1}\frac{(rt)^k}{k!},
\quad t\ge0.
\label{dnad24:erlang}
\end{equation}
One derivation counts arrivals of a rate-$r$ Poisson clock: the $d$th stage is
complete exactly when at least $d$ arrivals have occurred by $t$.
For zero stages, completion is immediate; $F_0(t)=1$.

The reference sums the upper Poisson tail when $rt<d$ to avoid subtracting a
number nearly one. For large $rt$, it subtracts the small lower tail.
Its declared numerical domain limits depth to 256; this is a reference budget,
not a physical limit on reactions.
\Listing{erlang_complete}

Now distinguish one completed computation from \emph{all} of $R$ independent
computations. Under independent chains, the latter probability is
\[
P(\max_jT_j\le t)=F_d(t)^R.
\]
Adding parallel jobs leaves the mean time of one chain unchanged but reduces the
probability that every job has completed by a fixed deadline.
If the system needs only one successful witness, it has a different event again.

\Plate{04-latency}{Four sequential stages at assigned rate
$r=.1\ \mathrm{s}^{-1}$. Curves are generated for completion of one chain and
all of ten or one hundred independent chains. Extra parallelism does not make
the all-complete event occur earlier.}{dnad24:latency}

\subsection{Where the chain model stops}
Real reactions can branch, compete for species, return to previous states, become
trapped or nucleate at different times. These mechanisms invalidate the equal-rate,
independent irreversible chain. Chapter 23's error model and this latency model
should not be multiplied unless their conditional relationship has been specified.
Waiting longer can improve completion and also increase leakage or degradation.
A deadline is therefore part of the scientific outcome, not an inconvenient
implementation detail.

\section{Transport has its own length scale}
For free Brownian diffusion in $d_{\mathrm{space}}\in\{1,2,3\}$ dimensions,
$\mathbb E[\|X(t)-X(0)\|^2]=2d_{\mathrm{space}}Dt$.
Setting the root-mean-square displacement equal to a length $a$ gives
\begin{equation}
t_{\mathrm{RMS}}=\frac{a^2}{2d_{\mathrm{space}}D}.
\label{dnad24:diffusion}
\end{equation}
Here $a$ is in metres and $D$ in $\mathrm{m}^2/\mathrm{s}$.
This is not a vessel mixing-time prediction or a first-passage time to a binding
site. Boundaries, geometry, stirring, crowding and reactive surfaces can change the
relevant transport problem.
\begin{center}\small\input{\Generated/transport.tex}\end{center}
The generated values use assigned $D=10^{-10}\ \mathrm{m}^2/\mathrm{s}$ in three
dimensions. Increasing travel distance by ten multiplies this RMS timescale by
one hundred. Scaling a drawing without changing transport assumptions does not
scale the device's kinetics faithfully.

\section{A correct molecule must enter the observed sample}
Readout is a physical channel, not a final line of pseudocode. Suppose the relevant
aliquot has a Poisson number of candidate witnesses with mean $\nu$.
Each is independently detectable with probability $\eta\in[0,1]$.
Poisson thinning makes the detected count Poisson with mean $\eta\nu$, hence
\begin{equation}
P(\text{at least one detected})=1-e^{-\eta\nu}.
\label{dnad24:detect}
\end{equation}
To keep the miss probability at most $\delta\in(0,1)$,
$\nu\ge-\log\delta/\eta$ when $\eta>0$.
At $\eta=0$, no finite witness inventory suffices. The implementation uses
\texttt{expm1} so that very small detection probabilities remain representable.
\Listing{detection_probability}
\begin{center}\small\input{\Generated/detection.tex}\end{center}

This event is not calibrated classification accuracy. It only describes capture
of at least one detectable molecule under a count model. A false-positive assay
can report a witness that was never there; Chapter 23 owns that separate
observation-channel inversion.

\Plate{05-readout}{Inventory, transport into an aliquot, detection and finite
report bandwidth are distinct gates. Counts describe expectations; the plate
does not promise that all molecules or all identities are observed.}
{dnad24:readout}

If a reader must communicate $Q$ bits over a serial channel with effective
bandwidth $b$ bits/s, its output alone requires at least $Q/b$ seconds under that
interface. A million bits at $10^5$ bits/s costs ten seconds before overhead.
This is not a measured sequencing rate. It illustrates why returning an entire
candidate set and returning one decision bit have different output costs.

\section{Energy claims require a thermodynamic contract}
Do not compare an ideal molecular transition to the wall-plug energy of a complete
electronic service. Energy accounting must include preparation, temperature
control, mixing, instrumentation, observation and downstream analysis within the
declared system boundary.

For an ideal quasistatic erasure of one unbiased classical bit at temperature
$T$, the familiar lower bound is $k_BT\log2$ \cite{Landauer1961}.
The reference calculates approximately $2.87\times10^{-21}$ J at 300 K.
To see the information premise, start with two equally likely logical states.
Their entropy is $S_{\mathrm{bit}}=-k_B\sum_{i=0}^1p_i\log p_i=k_B\log2$.
Resetting both to one known state makes that logical entropy zero. In an
isothermal process with no compensating entropy retained elsewhere, the
environment must gain at least $k_B\log2$, giving the ideal heat floor
$Q_{\mathrm{env}}\ge k_BT\log2$. The state probabilities and system boundary
are essential to the argument.
This is neither an empirical DNA energy cost nor a universal charge per hybridization
event. Biased states, retained correlations, reversibility and nonequilibrium
operation require the corresponding thermodynamic analysis. A useful engineering
measurement is actual energy per verified completed task, with failures and
auxiliary systems included.

\section{The Pareto problem: costs move rather than disappear}
A design $A$ dominates $B$ only relative to specified metrics and matched success
requirements: no controlled cost is worse and at least one is strictly better.
Otherwise, the designs may be incomparable.
A smaller species alphabet may require a larger seed; more copies may improve
sampling while increasing volume or readout; stronger redundancy may decrease
logical error while increasing area and serial depth.

\Plate{06-tradeoff}{A resource comparison freezes input, correctness event,
deadline and observation contract before placing designs on a Pareto frontier.
Numbers are not invented to rank real devices.}{dnad24:tradeoff}

For a molecular/electronic comparison, charge each system for input preparation,
representation, execution, output verification and recovery from failure.
Different devices may still be useful for different tasks. The goal is to discover
the workload where a complete design has an advantage, not to preserve an
unqualified story about exponential parallel speedup.

\section{Executable checks and falsifiers}
The chapter adds thirty-six local tests. Coverage inversion is checked at one,
eight, 1024 and $2^{60}$ species; exact small fixed-draw occupancy is compared
against independent exhaustive enumeration. A seeded Poisson-arrival Monte Carlo
checks a different preparation contract rather than accidentally validating the
fixed-draw law. A depth-two closed form independently verifies the Erlang routine.

Other checks enforce concentration/volume conservation, quadratic RMS transport
scaling, detection inversion, endpoints and invalid domains. These are tests of
declared mathematical models. They do not verify synthesis feasibility,
sequence-dependent rate constants, physical mixing or actual instrument sensitivity.

The main falsifiers are concrete: a zero-probability required species defeats
coverage; a depleted partner defeats constant-rate waiting; common resource
coupling defeats independent-chain completion; an insensitive assay defeats
detection; an uncounted preparation cost defeats a headline system comparison.

\begin{KeyIdeas}
\begin{itemize}\setlength{\itemsep}{2pt}
\item Compare designs only through a declared resource vector: sequences,
molecules, volume, concentration, length, time, readout and energy.
\item Poisson coverage: to include each of $M$ needed species with probability
$1-\delta$, the mean copies per species must satisfy
$\mu\ge-\log[1-(1-\delta)^{1/M}]$.
\item At fixed total concentration, more distinct species means fewer partners of
each kind, so bimolecular encounter times grow with diversity.
\item Parallelism shortens width, not depth: a chain of $d$ dependent steps has an
Erlang completion time.
\item Diffusion over distance $a$ takes about $a^2/(2dD)$; transport is a cost
when systems grow.
\item Landauer's bound is far below any laboratory overhead; a credible energy
claim states its boundary and includes synthesis, temperature control and
readout.
\end{itemize}
\end{KeyIdeas}

\section{Exercises and worked solutions}
\begin{enumerate}
\item Why does $2^n$ describe this search and not all DNA algorithms?
\textbf{Answer:} this architecture explicitly materializes every $n$-bit candidate;
another representation or search distribution can require different inventory.
\item Derive Equation~\ref{dnad24:coverage}.
\textbf{Answer:} each independent count is nonzero with probability $1-e^{-\mu}$;
multiply over $M$ species.
\item Solve the coverage requirement for $\mu$.
\textbf{Answer:} take the $M$th root of $(1-e^{-\mu})^M\ge1-\delta$ and then
the negative logarithm to obtain Equation~\ref{dnad24:mu}.
\item Explain the numerical use of \texttt{expm1}.
\textbf{Answer:} it accurately evaluates $e^x-1$ when $x$ is close to zero;
direct subtraction loses relative precision.
\item Derive the large-$M$ copy law.
\textbf{Answer:} with $a=-\log(1-\delta)>0$,
$1-e^{-a/M}=a/M+O(M^{-2})$, so $\mu=\log M-\log a+o(1)$.
\item Prove the union bound's sufficient requirement.
\textbf{Answer:} sum the $M$ missing probabilities $e^{-\mu}$ and solve
$Me^{-\mu}\le\delta$.
\item Why is a fixed total of $M\mu$ draws not the same model?
\textbf{Answer:} conditioning the total count removes independent species counts;
use the occupancy recurrence or another fixed-draw law.
\item Derive the next-draw transitions for $s$ represented species.
\textbf{Answer:} a represented identity is drawn with probability $s/M$;
a new identity with probability $(M-s)/M$.
\item If one required identity is absent, does amplification repair coverage?
\textbf{Answer:} no; descendants of present templates do not create the absent
identity under an identity-preserving amplification model.
\item Double total concentration with the same inventory. What changes?
\textbf{Answer:} modeled volume halves. Encounter rates may change; feasibility
and off-target reactions require separate analysis.
\item Derive Equation~\ref{dnad24:wait}.
\textbf{Answer:} equal species concentration is $c_{\mathrm{tot}}/S$;
the exponential mean is its reciprocal rate.
\item Compute the mean of four serial stages at $.1\ \mathrm{s}^{-1}$.
\textbf{Answer:} 40 s; this is a mean, not a completion guarantee at 40 s.
\item Derive the chain CDF from a Poisson clock.
\textbf{Answer:} completion requires at least $d$ arrivals; subtract the probability
of zero through $d-1$ arrivals.
\item Does running one hundred independent chains make all complete by the
one-chain median? \textbf{Answer:} probability is $(1/2)^{100}$, not $1/2$.
\item Scale RMS diffusion distance by ten.
\textbf{Answer:} the assigned-model timescale increases by one hundred, with $D$
and spatial dimension fixed.
\item How many mean witnesses are needed for .95 detection at $\eta=.1$?
\textbf{Answer:} $-\log(.05)/.1\approx29.9573$.
\item What is the readout lower bound for $10^6$ bits at $10^5$ bits/s?
\textbf{Answer:} ten seconds for the specified serial interface alone.
\item Why is $k_BT\log2$ not DNA energy per operation?
\textbf{Answer:} it is an ideal erasure bound under particular information and
thermodynamic premises, not a measured chemical-device energy budget.
\item Design a Pareto comparison.
\textbf{Rubric:} freeze task, representation, success event, deadline and observed
output; count species, copies, auxiliaries, synthesis, volume, latency, readout,
energy boundary and failed runs; report incomparable tradeoffs honestly.
\item Extend the coverage model to unequal abundances.
\textbf{Answer:} independent Poisson means $\mu_j$ give
$\prod_j(1-e^{-\mu_j})$; a zero mean makes complete coverage probability zero.
A useful extension optimizes inventory subject to measured abundance uncertainty.
\end{enumerate}

\section{Bridge}
Chapter 25 specializes the representation and readout account to DNA storage:
indices, payload constraints, strand dropout and decoding. Chapter 26 owns
laboratory workflows, controls and immutable evidence bundles. Keeping those
jobs distinct prevents resource theory from repeating a methods checklist.
The question changes from resource inventory to what information survives the
complete physical channel.
